% from ega4-10-fr.tex
\oldpage[IV-3]{107}
intégralement clos et de dimension $1$ et une $A$-algèbre finie~; c'est donc
un anneau de valuation discrète, et la proposition en résulte aussitôt dans
ce cas.

Passons au cas général~; soit $x$ le point fermé de
$X=\operatorname{Spec}(A)$, et posons $Y=X-\{x\}$~; on sait (10.5.9) que $Y$
est un préschéma de Jacobson. Montrons que cela entraîne que l'intersection
des idéaux premiers $\mathfrak p$ de $A$ tels que
$\dim(A/\mathfrak p)=1$ est \emph{réduite à $0$}~; la proposition en résultera
puisque, pour chacun de ces idéaux $\mathfrak p$, l'intersection des noyaux des
homomorphismes locaux de $A/\mathfrak p$ dans des anneaux de valuation
discrète $V$, faisant de $V$ une $(A/\mathfrak p)$-algèbre finie, est réduite
à $0$. Mais dire que l'intersection de ces idéaux premiers est réduite à $0$
signifie que l'ensemble de ces idéaux est dense dans $X$, ou encore dans $Y$
(puisque $Y$ est dense dans $X$), et cela résulte aussitôt de (10.5.9).

\subsection{Dimension dans les préschémas de Jacobson}
\label{subsection:IV.10.6-fr}

Les résultats de ce numéro précisent dans certains cas et généralisent des résultats
du \S~5.

\begin{proposition}[10.6.1]
\label{IV.10.6.1-fr}
Soit $S$ un préschéma localement noethérien, vérifiant en outre les conditions
suivantes~: $1^\circ$ $S$ est un préschéma de Jacobson~; $2^\circ$ pour tout
$s\in S$, $\mathcal O_{S,s}$ est universellement caténaire (5.6.2)~; $3^\circ$ toute
composante irréductible $S'$ de $S$ est équicodimensionnelle (autrement dit, pour
tout point fermé $s$ de $S'$ et tout sous-préschéma de $S$ ayant $S'$ pour espace
sous-jacent, on a $\dim(S')=\dim(\mathcal O_{S',s})$). On a alors les propriétés
suivantes~:
\begin{enumerate}[label=(\roman*)]
  \item Pour tout morphisme localement de type fini $g:X\to S$, $X$ vérifie les
  conditions $1^\circ$, $2^\circ$ et $3^\circ$ précédentes. En particulier, si $X$
  est équidimensionnel (par exemple si $X$ est irréductible), $X$ est
  biéquidimensionnel (autrement dit, $X$ est caténaire et pour tout point fermé
  $x$ de $X$, on a $\dim(\mathcal O_{X,x})=\dim(X)$ ($\mathbf 0$, 14.3.3)).
  \item Soient $X$, $Y$ deux $S$-préschémas localement de type fini sur $S$,
  $f:X\to Y$ un $S$-morphisme~; on suppose $X$ irréductible et $f$ dominant. Si
  $\xi$ (resp. $\eta$) est le point générique de $X$ (resp. $Y$) et
  $e=\dim(f^{-1}(\eta))=\deg.\operatorname{tr}_{k(\eta)}k(\xi)$, on a
  \begin{equation}
  \label{IV.10.6.1.1-fr}
  \tag{10.6.1.1}
  \dim(X)=\dim(Y)+e.
  \end{equation}
  \item Soient $X$, $Y$ deux $S$-préschémas localement de type fini sur $S$,
  $f:X\to Y$ un $S$-morphisme, $n$ un entier $\geq0$ tel que l'on ait
  $\dim(f^{-1}(y))\geq n$ (resp. $\dim(f^{-1}(y))\leq n$) pour tout $y\in Y$.
  Alors on a
  \begin{equation}
  \label{IV.10.6.1.2-fr}
  \tag{10.6.1.2}
  \dim(X)\geq\dim(Y)+n
  \end{equation}
  (resp.
  \begin{equation}
  \label{IV.10.6.1.3-fr}
  \tag{10.6.1.3}
  \dim(X)\leq\dim(Y)+n).
  \end{equation}
\end{enumerate}
\end{proposition}

(i) La propriété $1^\circ$ pour $X$ résulte de (10.4.7). Pour tout $x\in X$,
$\mathcal O_{X,x}$ est anneau local d'une $\mathcal O_{S,g(x)}$-algèbre de type fini en
un idéal premier, et l'homomorphisme
$\mathcal O_{S,g(x)}\to\mathcal O_{X,x}$

